\section{Real parameter tests and the background contraction}
\label{sec:background-ray-estimates}
The small Lipschitz constant used in the background history does not
follow from a small bound on the values alone. It follows from the
paid bounds on finite linear combinations along each real parameter
line. We prove the conversion explicitly. Only real configurations
are used in these tests.

\subsection{Recovering a convergent series from finite tests}
Let $E$ be a real Banach space, $p>1$, $C\ge0$, and
$f:[-1,1]\to E$. Suppose that, for every finite list
$t_1,\ldots,t_k\in[-1,1]$ and real weights $a_1,\ldots,a_k$,
\begin{equation}\label{eq:ray-dual-test}
 \left\|\sum_{j=1}^k a_j f(t_j)\right\|
 \le C\sup_{n\ge0}\left|\sum_{j=1}^k a_j(t_j/p)^n\right|.
\end{equation}
No continuity of $f$ is assumed here. There are vectors $b_n\in E$
with $\|b_n\|\le C$ such that
\begin{equation}\label{eq:ray-recovered-series}
                    f(t)=\sum_{n=0}^\infty(t/p)^n b_n
                    \qquad(-1\le t\le1).
\end{equation}

To prove this, use the sup norm on sequences and let
$v(t)=((t/p)^n)_{n\ge0}$. Its coordinates tend to zero. Define,
on the span of these sequences,
\[
                 T\left(\sum_j a_jv(t_j)\right)=\sum_j a_jf(t_j).
\]
If the sequence on the left is zero, \eqref{eq:ray-dual-test}
makes the proposed value zero. Thus $T$ is well defined and linear,
and its norm is at most $C$. It extends uniquely to the closed
span, because $E$ is complete.

Every coordinate vector $e_n$ lies in that closed span. For $n=0$,
use $v(0)=e_0$. Inductively assume this for $e_0,\ldots,e_{n-1}$.
For $0<t\le1$, put $r=t/p$. The vector
\[
       r^{-n}\left(v(t)-\sum_{j<n}r^je_j\right)
\]
has its first $n$ coordinates zero, its $n$th coordinate one,
and subsequent coordinates $r,r^2,\ldots$. Its sup distance
from $e_n$ is $r\to0$ as $t\downarrow0$. This proves the claim.
Set $b_n=T e_n$, so $\|b_n\|\le C$. The geometric series
$v(t)=\sum_n(t/p)^ne_n$ converges in the sup norm. Apply $T$
to it to obtain \eqref{eq:ray-recovered-series}.

For $1\le a<p$, the same identity can be written
$f(t)=\sum_n(t/a)^nc_n$ with $c_n=(a/p)^nb_n$ and
\begin{equation}\label{eq:ray-radius-loss}
                        \sum_n\|c_n\|\le\frac{C}{1-a/p}.
\end{equation}
This is the radius loss used when passing from finite tests to
absolutely summable analytic coefficients. In particular, taking
$a=p/2$ costs at most a factor two when $p\ge2$. The argument
also applies to maps into the completion of a seminormed space;
no completeness of its original representatives is needed.

Two consequences will be used. First, a test with one point and
weight one gives $\|f(t)\|\le C$. Second, subtracting the series
at zero gives
\begin{equation}\label{eq:ray-increment}
                 \|f(t)-f(0)\|\le C\frac{|t|}{p-|t|}.
\end{equation}
Both conclusions hold even if the initial pointwise construction
of $f$ involved arbitrary choices: the finite tests themselves
force the displayed series and hence continuity on $[-1,1]$.

\subsection{Applying the tests to the actual activity coordinates}
At one adjacent pair of scales, retain the actual common dictionary
of regular and forced arrays. Its norm is the maximum of the
normalized regular and forced budgets, with all the fixed symmetry
and support constraints retained. Finite linear combinations remain
in this vector space. Suppose a line of allowed inputs has normalized
coupling $\lambda+t\mu$ and activity array $x+ty$, $|t|\le1$.
Let $F(t)$ be its selected output array in the completion and
$Q(t)$ its scalar coupling correction.

The exact finite-test receipts used here are the following. For
any finite weights $a_j$ and points $t_j$, write
\[
 \eta_r=\sup_{n\ge0}\left|\sum_j a_j(t_j/r)^n\right|.
\]
The combined regular output array has normalized budget at most
$d\eta_{q/2}$, as does the combined forced output array. The
scalar correction satisfies
\[
                   \left|\sum_j a_jQ(t_j)\right|
                          \le f_m\eta_q.
\]
The forced statement follows from its single common summable
coefficient majorant; the regular statement is a bound on the
actual mixed arrays, with their derivative prices, not just on
values of their gas. The definition of the array norm and the
isometric embedding into its completion consequently give
\begin{equation}\label{eq:background-ray-tests}
 \left\|\sum_j a_jF(t_j)\right\|\le d\eta_{q/2},\qquad
                 \left|\sum_j a_jQ(t_j)\right|\le f_m\eta_q.
\end{equation}
These are precisely the two instances of
\eqref{eq:ray-dual-test} needed below. Padding the dictionaries
with zero entries, or using another representative of the same
array, preserves them.

We spell out why the admissible lines are the ones used in the
finite-test construction. Their activity values are affine in $t$
on each fixed support. On a healthy configuration the regular mask
is one and the value is the corresponding affine combination;
on an unhealthy configuration the common mask is zero and the
same equality holds. Forced values are affine without a derivative
condition. The inverse coupling along the line is
$\widehat\beta\{w(\lambda+t\mu)\}^2$, with the same reference
$\widehat\beta$ and $w=\widehat\beta^{-1/32}$. Thus it is exactly
the prescribed quadratic coupling line, not an independent choice
of complexified gauge fields. The protected supports, normalization,
path data and transition are the same for all $t$.
The analytic estimates that pay \eqref{eq:background-ray-tests}
are retained as quantitative inputs here; the deduction of the
contraction from those estimates is complete below.

\subsection{The full-band Lipschitz estimates}
Take $q\ge4$, $d,f_m\ge0$, $2d\le\gamma=1/128$, and $s_m=2f_m$.
Assume $w\le1/128$, as purchased by the coupling threshold.
Consider two inputs $z=(\lambda,x)$ and
$z'=(\lambda',x')$ with $\|x\|,\|x'\|\le1/8$ and
$17/20\le w\lambda,w\lambda'\le23/20$. Use the product
norm
\[
                 h=\|z'-z\|=\max\{|\lambda'-\lambda|,\|x'-x\|\}.
\]
The constant-line tests give $\|F(z)\|,\|F(z')\|\le d$
and $|Q(z)|,|Q(z')|\le f_m$. If $h>1/4$, the triangle
inequality therefore gives
\[
 \|F(z')-F(z)\|\le2d\le8\gamma h,
 \qquad |Q(z')-Q(z)|\le2f_m\le8f_m h=4s_mh.
\]
For $h=0$, no equality of raw representatives is needed.
For any $0<\varepsilon\le1$, take the line
$z+t(z'-z)/\varepsilon$, $|t|\le1$. Its coupling direction is
zero and its activity direction has seminorm zero. Every point
therefore has the same activity norm as $z$, at most $1/8$,
so the line is admissible. Its value at $t=\varepsilon$ is $z'$.
Equation~\eqref{eq:ray-increment} bounds the output differences
by $d\varepsilon/(q/2-\varepsilon)$ and
$f_m\varepsilon/(q-\varepsilon)$, respectively. Letting
$\varepsilon\downarrow0$ proves equality of the output classes
and of the scalar correction. This also proves that the map
respects the seminorm kernel; it does not identify regular
representatives on unhealthy configurations.

For $0<h\le1/4$, set $v=(z'-z)/(4h)$. Its norm is $1/4$,
and the full line $z+tv$, $|t|\le1$, is allowed. Its activity
norm is at most $1/8+1/4<1$. Its normalized coupling lies in
the broader interval
\[
                   33/40\le w(\lambda+tv_1)\le47/40,
\]
since $w|tv_1|\le1/512<1/40$, the margin between the narrower
and broader intervals. The prescribed coupling-line construction
therefore supplies \eqref{eq:background-ray-tests} on this line.
Apply \eqref{eq:ray-increment} at $t=4h\le1$ to get
\[
 \|F(z')-F(z)\|\le\frac{4d}{q/2-1}h\le4dh\le8\gamma h,
 \qquad
 |Q(z')-Q(z)|\le\frac{4f_m}{q-1}h\le8f_mh.
\]
This proves exactly \eqref{eq:coupled-activity-step} and
\eqref{eq:coupled-cost-step}. The output bounds also follow from
the constant-line tests, since $d\le1/256<1/8$ and $f_m=s_m/2$.

The use of one-dimensional series here is sufficient for these
Lipschitz estimates. The same real finite-test method also produces
multilinear Taylor coefficients by polarization when a full analytic
parameterization is required. That additional conclusion is not
needed in the convergence argument of
Section~\ref{sec:coupled-background-convergence}.
